\subsubsection*{Свойства формул обновления}

\begin{idea}[title={Откуда этот материал}]
На лекции формулы ДФП и БФГШ были выписаны без обоснования. Ниже доказано, что
обе удовлетворяют секущему условию и сохраняют положительную определённость —
то есть что направление $p_k = -A_k \nabla f(x_k)$ остаётся направлением
спуска.
\end{idea}

Перейдём к стандартным обозначениям: $s_k = x_{k+1} - x_k$,
$y_k = \nabla f(x_{k+1}) - \nabla f(x_k)$ (в обозначениях лекции
$s_k = \Delta x_k$, $y_k = -\Delta\omega_k$). \textbf{Секущее
условие} — это требование
\[
  A_{k+1}\, y_k = s_k ,
\]
а \textbf{условие кривизны} — $s_k\T y_k > 0$.

\begin{theorem}[ДФП: секущее условие]\label{thm:dfp-secant}
Пусть
\[
  A_{k+1} = A_k + \frac{s_k s_k\T}{s_k\T y_k}
            - \frac{A_k y_k y_k\T A_k}{y_k\T A_k y_k} .
\]
Тогда $A_{k+1} y_k = s_k$.
\end{theorem}

\begin{proof}
Умножим формулу справа на $y_k$:
\[
  A_{k+1} y_k = A_k y_k + \frac{s_k \,(s_k\T y_k)}{s_k\T y_k}
    - \frac{A_k y_k\,(y_k\T A_k y_k)}{y_k\T A_k y_k} .
\]
Величины $s_k\T y_k$ и $y_k\T A_k y_k$ — числа, поэтому второе слагаемое равно
$s_k$, а третье равно $A_k y_k$. Следовательно
$A_{k+1} y_k = A_k y_k + s_k - A_k y_k = s_k$.
\end{proof}

\begin{theorem}[ДФП: положительная определённость]\label{thm:dfp-pd}
Пусть $A_k \succ 0$ и $s_k\T y_k > 0$. Тогда $A_{k+1} \succ 0$.
\end{theorem}

\begin{proof}
Для произвольного $z \ne 0$
\[
  z\T A_{k+1} z = z\T A_k z
    - \frac{\big(z\T A_k y_k\big)^2}{y_k\T A_k y_k}
    + \frac{\big(z\T s_k\big)^2}{s_k\T y_k} .
\]
Так как $A_k \succ 0$, выражение
$\langle u, v\rangle_{A_k} = u\T A_k v$ — скалярное произведение, и по
неравенству Коши--Буняковского
\[
  \big(z\T A_k y_k\big)^2 \le \big(z\T A_k z\big)\big(y_k\T A_k y_k\big)
  \qquad\Longrightarrow\qquad
  z\T A_k z - \frac{\big(z\T A_k y_k\big)^2}{y_k\T A_k y_k} \ \ge\ 0 .
\]
Второе слагаемое неотрицательно, так как $s_k\T y_k > 0$. Значит
$z\T A_{k+1} z \ge 0$.

Покажем, что равенство невозможно при $z \ne 0$. Если $z\T A_{k+1}z = 0$, то
оба неотрицательных слагаемых нулевые. Из обращения в равенство неравенства
Коши следует, что $z$ пропорционален $y_k$: $z = c\,y_k$. Из второго
слагаемого $\big(z\T s_k\big)^2 = 0$, то есть $z\T s_k = 0$, откуда
$c\,y_k\T s_k = 0$. Но $y_k\T s_k = s_k\T y_k > 0$, значит $c = 0$ и $z = 0$ —
противоречие. Следовательно $z\T A_{k+1} z > 0$ при всех $z \ne 0$.
\end{proof}

\begin{theorem}[БФГШ: секущее условие]\label{thm:bfgs-secant}
Пусть
\[
  A_{k+1} = \big(E - \rho_k s_k y_k\T\big) A_k \big(E - \rho_k y_k s_k\T\big)
            + \rho_k s_k s_k\T ,
  \qquad \rho_k = \frac{1}{y_k\T s_k} .
\]
Тогда $A_{k+1} y_k = s_k$.
\end{theorem}

\begin{proof}
Рассмотрим правый множитель:
\[
  \big(E - \rho_k y_k s_k\T\big) y_k = y_k - \rho_k y_k \big(s_k\T y_k\big)
  = y_k - y_k = 0 ,
\]
так как $\rho_k \big(s_k\T y_k\big) = 1$. Значит первое слагаемое обращается в
нуль, и остаётся
\[
  \rho_k s_k s_k\T y_k = \frac{1}{y_k\T s_k}\, s_k \big(s_k\T y_k\big) = s_k .
  \qedhere
\]
\end{proof}

\begin{theorem}[БФГШ: положительная определённость]\label{thm:bfgs-pd}
Пусть $A_k \succ 0$ и $y_k\T s_k > 0$. Тогда $A_{k+1} \succ 0$.
\end{theorem}

\begin{proof}
Обозначим $w = \big(E - \rho_k y_k s_k\T\big) z$; тогда
$z\T\big(E - \rho_k s_k y_k\T\big) = w\T$, и для произвольного $z$
\[
  z\T A_{k+1} z = w\T A_k w + \rho_k \big(z\T s_k\big)^2 .
\]
Оба слагаемых неотрицательны ($A_k \succ 0$ и $\rho_k > 0$), значит
$z\T A_{k+1}z \ge 0$.

Пусть $z\T A_{k+1}z = 0$. Тогда $w\T A_k w = 0$ и $\big(z\T s_k\big)^2 = 0$.
Из $A_k \succ 0$ следует $w = 0$; из второго равенства $s_k\T z = 0$. Но
\[
  w = z - \rho_k y_k \big(s_k\T z\big) = z ,
\]
и, поскольку $w = 0$, получаем $z = 0$. Значит $z\T A_{k+1}z > 0$ при
$z \ne 0$.
\end{proof}

\begin{idea}
Условие кривизны $s_k\T y_k > 0$ выполняется автоматически для сильно выпуклой
$f$ и при точном одномерном поиске; в общем случае его обеспечивают условиями
Вульфа на шаг. Если же $s_k\T y_k \le 0$, обе формулы могут потерять
положительную определённость — тогда обновление пропускают или делают сброс
$A_{k+1} = E$ (шаг~9 алгоритма §\ref{subsec:quasi-newton-anchor}).

БФГШ устойчивее ДФП именно из-за своей формы: обновление записано как
«сопряжение» $W A_k W\T$ плюс слагаемое ранга один, и положительная
определённость в нём не может потеряться из-за накопления ошибок округления,
тогда как в ДФП она держится на разности двух слагаемых.
\end{idea}
